\subsubsection{测试段二}

本节测试交换图, 脚注, 超链接与嵌套列表.

脚注测试: 这是正文文字\footnote{这是脚注内容.}

超链接测试: 链接到 \href{https://www.example.com}{示例网站}.

嵌套列表测试:
\begin{enumerate}
    \item 第一层第一项.
    \item 第一层第二项.
    \begin{enumerate}
        \item 第二层第一项.
        \item 第二层第二项.
    \end{enumerate}
    \item 第一层第三项.
\end{enumerate}

多标签引用测试: \cref{thm:test,def:test}, \cref{eq:cases-1,eq:matrix-1}.

\medskip

下面测试 tikz-cd 交换图.

张量积的泛性质 (图~\ref{cd:tensor}):

\begin{figure}[htbp]
    \centering
    \begin{tikzcd}
        M \times N \arrow[r, "B"] \arrow[d, "\otimes"'] & P \\
        M \otimes N \arrow[ru, "\tilde{B}"', dashed]
    \end{tikzcd}
    \caption{张量积的泛性质}
    \label{cd:tensor}
\end{figure}

五引理 (图~\ref{cd:five}):

\begin{figure}[htbp]
    \centering
    \begin{tikzcd}
        A_1 \arrow[r] \arrow[d, "\alpha_1"] & A_2 \arrow[r] \arrow[d, "\alpha_2"] & A_3 \arrow[r] \arrow[d, "\alpha_3"] & A_4 \arrow[r] \arrow[d, "\alpha_4"] & A_5 \arrow[d, "\alpha_5"] \\
        B_1 \arrow[r] & B_2 \arrow[r] & B_3 \arrow[r] & B_4 \arrow[r] & B_5
    \end{tikzcd}
    \caption{五引理}
    \label{cd:five}
\end{figure}

交换图引用测试: \autoref{cd:tensor}, \autoref{cd:five}.
